% ==================================================================
% id:       prob-012
% topic:    מבוא להסתברות
% subtopic: תכונות של פונקציית הסתברות
% solution: full
% ==================================================================
יהיו $A,B,C$ שלושה מאורעות זרים בזוגות, כך ש-$\PP(A)=\PP(B)=\PP(C)=p$. מה ההסתברות שיקרה בדיוק אחד מהם?
\begin{solbox}
המאורעות זרים, ולכן לא ייתכן שיקרו שניים מהם יחד. לכן „בדיוק אחד” הוא בדיוק „לפחות אחד”, כלומר $A\cup B\cup C$. כיוון שהם זרים בזוגות,
\[
\PP(A\cup B\cup C)=\PP(A)+\PP(B)+\PP(C)=\ans{3p}.
\]
\textbf{הערה.} מכאן גם ש-$3p\le1$, כלומר $p\le\frac13$: לא קיימים שלושה מאורעות זרים בזוגות, שההסתברות של כל אחד מהם היא, למשל, $\frac12$.
\end{solbox}
